% !TeX program = lualatex
% !TeX encoding = utf8
% !TeX spellcheck = uk_UA

% chktex-file 3
% chktex-file 21
% chktex-file 25
% chktex-file 46

\clearpage

%% --------------------------------------------------------
\simheading\label{sec:Simulations}\hypertarget{Simulations}{}
%% --------------------------------------------------------

Нижче зібрано всі інтерактивні симуляції підручника в~порядку розділів. Кожна відкривається в~браузері
без встановлення додаткових програм (посилання дієві в~електронній версії книжки); тут же вказано розділ, до якого належить
симуляція, і~коротко описано, що в~ній можна дослідити.

\begin{description}
	\item[Розділ~\ref{SteadyEfield}.]
	\item[] 
		\begin{itemize}
			\item \simlink{field_lines_equipotentials.html}{Силові лінії та~еквіпотенціалі системи точкових зарядів}~---
				картина поля й~еквіпотенціальних ліній для довільного набору зарядів.
		\end{itemize}

	\item[Розділ~\ref{FieldInDielectrics}.]
		\begin{itemize}
			\item \simlink{Polar_dielectric.html}{Полярний діелектрик}~---
				орієнтація диполів у~полі за~Ланжевеном--Дебаєм: залежність поляризації від поля й~температури.
		\end{itemize}

	\item[Розділ~\ref{SteadyCurrent}.]
		\begin{itemize}
			\item \simlink{rc_lamp_transients.html}{Конденсатор і~лампа розжарювання}~---
				зарядка й~розрядка конденсатора через лампу; опір нитки залежить від температури, тому криві відрізняються від експонент.
		\end{itemize}

	\item[Розділ~\ref{MagFieldMedia}.]
		\begin{itemize}
			\item \simlink{ferromagnetics.html}{Феромагнетик}~---
				перемагнічування, рух доменних стінок і~петля гістерезису.
		\end{itemize}

	\item[Розділ~\ref{ElMagInduction}.]
		\begin{itemize}
			\item \simlink{faraday_magnet_coil.html}{Дослід Фарадея: рух магніту або котушки}~---
				ЕРС залежить лише від відносного руху; магніт можна провести крізь котушку й~побачити зміну знака ЕРС.
			\item \simlink{faraday_switch_coils.html}{Дослід Фарадея: вмикання і~вимикання струму}~---
				взаємоіндукція двох котушок при замиканні й~розмиканні ключа.
			\item \simlink{betatron.html}{Бетатрон}~---
				наростання поля, розгін електрона, бетатронні коливання й~вивід на мішень.
			\item \simlink{rl_lamp_transients.html}{Котушка і~лампа розжарювання}~---
				встановлення й~спадання струму в~колі з~індуктивністю та~нелінійним опором.
		\end{itemize}

	\item[Розділ~\ref{Radiation}.]
		\begin{itemize}
			\item \simlink{thompson_radiation_model_acceleration.html}{Модель Томсона: прискорення заряду}~---
				імпульс випромінювання як збурення силових ліній, що поширюється від заряду.
			\item \simlink{thompson_radiation_model_decceleration.html}{Модель Томсона: гальмування заряду}~---
				те саме для різкої зупинки заряду.
			\item \simlink{hertz_dipole_radiation.html}{Диполь Герца}~---
				відрив силових ліній від диполя й~випромінювання.
			\item \simlink{Synchrotron_radiation.html}{Синхротронне випромінювання}~---
				конус випромінювання шириною $1/\gamma$ і~спалах у~спостерігача.
			\item \simlink{synchrotron_radiation_v2.html}{Синхротронне випромінювання: конус і~фронти хвиль}~---
				друга версія, що показує конус і~фронти хвиль.
			\item \simlink{cherenkov_radiation.html}{Черенковське випромінювання}~---
				конус хвиль від частинки, що рухається швидше за~фазову швидкість світла в~середовищі.
		\end{itemize}

	\item[Розділ~\ref{ElMagOscillations}.]
		\begin{itemize}
			\item \simlink{rlc_series_resonance.html}{Резонанс напруг}~---
				послідовний $RLC$-контур, осцилограми, резонансні криві та~вплив осердя в~котушці.
			\item \simlink{rlc_parallel_resonance.html}{Резонанс струмів}~---
				паралельний контур з~осердям у~котушці та~послідовним опором генератора.
		\end{itemize}

	\item[Розділ~\ref{Telegraph}.]
		\begin{itemize}
			\item \simlink{transmission_line.html}{Телеграфна лінія}~---
				падаюча й~відбита хвилі в~довгій лінії.
		\end{itemize}
\end{description}
